\documentclass[10pt,a4paper]{amsart}
\usepackage[margin=19mm]{geometry}
\usepackage{amsmath,amssymb,mathtools}
\usepackage[scr=rsfs,cal=euler]{mathalfa}
\usepackage{xfrac}
\usepackage[unicode,hidelinks,bookmarks=false]{hyperref}
\newcommand{\ie}{i.e.\ }
\newcommand{\cf}{cf.\ }
\newcommand{\resp}{respectively\ }
\newcommand{\Subsection}{\subsection}
\providecommand{\nobreakdash}{\nobreak\hbox{-}}
\setlength{\emergencystretch}{2em}
\multlinegap0pt
\usepackage{bm}
\usepackage{bbm}
\usepackage{dsfont}
\usepackage{xspace}
\usepackage{cancel}
\usepackage{mathtools}
\usepackage{amsthm}
\newtheorem{theorem}{Theorem}
\newtheorem{lemma}[theorem]{Lemma}
\newcommand{\C}{\mathbb C}
\newcommand{\N}{\mathbb N}
\newcommand{\Wr}{\operatorname{Wr}}
\newcommand{\cW}{\mathcal W}
\newcommand{\ord}{\operatorname{ord}}
\title[The frame set of the first Hermite function]{The frame set of the first Hermite function}
\author{Markus Faulhuber, Philipp Petersen}
\date{Source: arXiv:2609.02610v2 (CC BY 4.0)}
\begin{document}
\maketitle

\noindent\emph{Source and licence.} The frame set of the first Hermite function. Version arXiv:2609.02610v2 (CC BY 4.0), distributed under the Creative Commons Attribution 4.0 International licence, as recorded in the arXiv licence metadata for this exact version and shown on the abstract page https://arxiv.org/abs/2609.02610v2. Full rights notice: licenses/arxiv-2609.02610-CC-BY-4.0.txt.\par\smallskip

\noindent\textit{Editorial scope.} Counted unit: the Lemma whose source label is \texttt{lem:wronskian-amplification} in the article (source0.tex lines 549--573, source label \texttt{lem:wronskian-amplification}), delivered together with the article's own complete original proof of it (source0.tex lines 575--643, one \texttt{proof} environment of 69 lines). No mathematics is written, repaired or re-derived: statement and proof are the article's own text line for line. Required context retained uncounted: none. Every definition, display and label of those lines is retained unchanged. Omitted from this package and not redistributed: 1--548, 574--574, 644--1054. Nothing inside a retained excerpt is omitted or altered: every retained body is delivered exactly as the article's lines are written, so no omission receipt of any kind accompanies this package. Citations: the retained excerpts carry no citation command at all, so no bibliography entry of the article is transcribed and no reference is redistributed. Reference adaptation: none was needed and none was made; every reference inside the delivered unit resolves inside it, so no unresolved reference or citation is produced. Printed slips kept literal: no printed word, symbol, hypothesis or number of the counted statement or of its proof was altered. Layout and class: the article's own class is \texttt{amsart} with packages including geometry, amsmath, amssymb, amsthm, mathtools, natbib, doi, xcolor, hyperref, enumerate; the wrapper is the lane's pinned \texttt{amsart} editorial wrapper at 10pt on A4, which restates the theorem-like environments the retained text needs and adds the article's own macro definitions that the retained text needs (\texttt{\textbackslash C}, \texttt{\textbackslash N}, \texttt{\textbackslash Wr}, \texttt{\textbackslash cW}, \texttt{\textbackslash ord}), with extra wrapper packages (bm, bbm, dsfont, xspace, cancel, mathtools). The delivered numbering is the unit's own. Assets: no retained span contains a figure environment, a graphics-inclusion command, a PDF-page inclusion, a table or a TikZ picture; no image file is copied into this package, the package declares no asset dependency and no image file is redistributed with it. Licence: version arXiv:2609.02610v2 of this article is distributed under the Creative Commons Attribution 4.0 International licence, as recorded in the arXiv licence metadata for this exact version and shown on the abstract page https://arxiv.org/abs/2609.02610v2. A full rights notice is delivered at licenses/arxiv-2609.02610-CC-BY-4.0.txt. Characters: every retained excerpt is pure ASCII, so the delivered engine, XeLaTeX with its default Latin Modern text font, reports no missing character.
\begin{lemma}
    \label{lem:wronskian-amplification}
    Let $N \in \N$, $N \geq 2$, let $U\subset\C$ be a connected open neighborhood of $\lambda\in\C$,
    and let $f_0,\ldots,f_{N-1}$ be holomorphic on $U$. Denote the Wronskian by
    \begin{equation}\label{eq:Wr-amplification-W}
        \cW(z)
        := \Wr(f_0,\ldots,f_{N-1})(z)
        = \det\bigl(f_j^{(r)}(z)\bigr)_{0\leq r,j<N},
        \qquad \text{ for } z\in U.
    \end{equation}
    Suppose that
    \begin{equation}\label{eq:Wr-amplification-hypothesis}
        f_j'(\lambda)=0,
        \qquad \text{ for all } j=0,\ldots,N-1.
    \end{equation}
    Then
    \begin{equation}\label{eq:Wr-derivative-vanishing}
        \cW^{(m)}(\lambda)=0,
        \qquad \text{ for all } m=0,\ldots,N-2.
    \end{equation}
    Consequently, if $\cW\not\equiv 0$, then
    \begin{equation}\label{eq:Wr-order}
        \ord_{\lambda}\cW\geq N-1.
    \end{equation}
\end{lemma}

\begin{proof}
    For nonnegative integers $r_0,\ldots,r_{N-1}$ and $z\in U$, set
    \begin{equation}\label{eq:D-determinant}
        D(r_0,\ldots,r_{N-1};z)
        :=
        \det\bigl(f_j^{(r_k)}(z)\bigr)_{0\leq k,j<N}.
    \end{equation}
    Thus
    \begin{equation}\label{eq:W-as-D}
        \cW(z)=D(0,1,\ldots,N-1;z).
    \end{equation}
    The point of the underlying argument now is that, after fewer than $N-1$ differentiations, every determinant arising from the Leibniz rule either has two equal rows or contains the row of first derivatives, which vanishes at $\lambda$.
    
    We start with repeated differentiation of \eqref{eq:W-as-D}, which, together with the multilinearity of the determinant in its rows, gives
    \begin{align}\label{eq:expansionOfWm}
        \cW^{(m)}(z)
        =
        \sum_{\substack{\alpha\in\N_0^N\\ |\alpha|=m}}
        \binom{m}{\alpha}
        D\bigl(
            \alpha_0,
            1+\alpha_1,
            \ldots,
            N-1+\alpha_{N-1};
            z
        \bigr), \quad \text{ for all } z \in U,
    \end{align}
    where
    \begin{equation}
        |\alpha|:=\alpha_0+\cdots+\alpha_{N-1},
        \qquad
        \binom{m}{\alpha}
        :=
        \frac{m!}{\alpha_0!\cdots\alpha_{N-1}!}.
    \end{equation}
    
    Fix $m\in\{0,\ldots,N-2\}$ and consider one term in the sum \eqref{eq:expansionOfWm}. Write
    \begin{equation}\label{eq:row-indices}
        r_k:=k+\alpha_k,
        \qquad \text{ for all } k=0,\ldots,N-1.
    \end{equation}
    If two of the integers $r_0,\ldots,r_{N-1}$ coincide, then the corresponding determinant in \eqref{eq:expansionOfWm}, defined in \eqref{eq:D-determinant}, has two equal rows and therefore vanishes.
    
    Suppose instead that $r_0,\ldots,r_{N-1}$ are pairwise distinct. If none of them were equal to $1$, then they would be $N$ distinct nonnegative integers avoiding $1$. Their sum would therefore satisfy
    \begin{equation}\label{eq:row-index-lower-bound}
        \sum_{k=0}^{N-1}r_k
        \geq
        0+2+3+\cdots+N
        =
        \frac{N(N-1)}{2}+N-1.
    \end{equation}
    On the other hand,
    \begin{equation}\label{eq:row-index-upper-bound}
        \sum_{k=0}^{N-1}r_k
        =
        \sum_{k=0}^{N-1}k+\sum_{k=0}^{N-1}\alpha_k
        =
        \frac{N(N-1)}{2}+m
        \leq
        \frac{N(N-1)}{2}+N-2,
    \end{equation}
    which contradicts \eqref{eq:row-index-lower-bound}. Hence, one of the row indices $r_k$ defined in \eqref{eq:row-indices} must equal $1$.
    
    The row whose index equals $1$ is
    \begin{equation}\label{eq:first-derivative-row}
        \bigl(f_0'(\lambda),\ldots,f_{N-1}'(\lambda)\bigr).
    \end{equation}
    By \eqref{eq:Wr-amplification-hypothesis}, we have that \eqref{eq:first-derivative-row} is zero. Thus every term in the expansion \eqref{eq:expansionOfWm} of $\cW^{(m)}(\lambda)$ vanishes, proving \eqref{eq:Wr-derivative-vanishing}. If $W\not\equiv0$, the vanishing of these derivatives gives \eqref{eq:Wr-order}, i.e., $\ord_\lambda \cW \geq N-1$.
\end{proof}

\end{document}
